\textbf{1. Kiến thức cần nhớ}
\begin{itemize}
\item \textbf{Hình chữ nhật} là tứ giác có bốn góc vuông ($\widehat{A} = \widehat{B} = \widehat{C} = \widehat{D} = 90^\circ$).
\begin{itemize}
\item Hình chữ nhật cũng là một hình bình hành, cũng là một hình thang cân.
\end{itemize}
\item \textbf{Tính chất}: Hình chữ nhật có đầy đủ tính chất của hình bình hành và hình thang cân:
\begin{itemize}
\item Các cạnh đối song song và bằng nhau.
\item \textbf{Hai đường chéo bằng nhau và cắt nhau tại trung điểm của mỗi đường}:
80669\boxed{AC = BD\quad \text{và}\quad OA = OB = OC = OD}80669
\end{itemize}
\item \textbf{Áp dụng vào tam giác vuông}:
\begin{itemize}
\item Trong tam giác vuông, \textbf{đường trung tuyến ứng với cạnh huyền bằng nửa cạnh huyền}:
80669\boxed{AM = \dfrac{1}{2}BC\quad (\triangle ABC\text{ vuông tại } A,\ M\text{ là trung điểm } BC)}80669
\item Đảo lại: Nếu một tam giác có đường trung tuyến ứng với một cạnh bằng nửa cạnh đó thì tam giác đó là \textbf{tam giác vuông}.
\end{itemize}
\item \textbf{Dấu hiệu nhận biết hình chữ nhật}:
\begin{enumerate}
\item Tứ giác có ba góc vuông là hình chữ nhật.
\item Hình thang cân có một góc vuông là hình chữ nhật.
\item Hình bình hành có một góc vuông là hình chữ nhật.
\item Hình bình hành có hai đường chéo bằng nhau là hình chữ nhật.
\end{enumerate}
\end{itemize}

\textbf{2. Phương pháp giải các dạng}

%% DANG 8H3?13-1
\textbf{Dạng 1. Tính chất hình chữ nhật, áp dụng vào tam giác vuông}
\begin{enumerate}
\item Sử dụng tính chất $ đoạn thẳng từ tâm đối xứng đến $ đỉnh bằng nhau:  = OB = OC = OD$ để chứng minh các tam giác cân hoặc tính độ dài.
\item Trong tam giác vuông, khi biết cạnh huyền tính được trung tuyến ứng với cạnh huyền và ngược lại.
\item Chứng minh một tam giác là tam giác vuông bằng cách chứng minh đường trung tuyến bằng nửa cạnh tương ứng.
\end{enumerate}

%% DANG 8H3?13-2
\textbf{Dạng 2. Dấu hiệu nhận biết hình chữ nhật}
\begin{enumerate}
\item Cách 1: Chứng minh tứ giác có $ góc vuông.
\item Cách 2: Chứng minh tứ giác là hình bình hành, sau đó chỉ ra:
\begin{itemize}
\item Có một góc vuông.
\item Hoặc có hai đường chéo bằng nhau ( = BD$).
\end{itemize}
\end{enumerate}

%% DANG 8H3?13-3
\textbf{Dạng 3. Tính độ dài, bài toán thực tế}
\begin{enumerate}
\item Áp dụng định lí Pythagore trong tam giác vuông tạo bởi hai cạnh kề và đường chéo của hình chữ nhật:
80669d^2 = a^2 + b^2 \implies \boxed{d = \sqrt{a^2 + b^2}}80669
\item Bài toán thực tế: Kiểm tra góc vuông khi thi công (thợ xây đo hai đường chéo bằng nhau để kiểm tra góc vuông), tính khoảng cách, kích thước màn hình tivi (đường chéo).
\end{enumerate}
